\documentclass[11pt]{article}
\usepackage[margin=1in]{geometry}
\usepackage[T1]{fontenc}
\usepackage{lmodern}
\usepackage{amsmath,amssymb}
\usepackage{hyperref}

\title{Literature Status for arXiv:math/0210228: \(L_p\) and Partition--Weight Norms}
\author{Run fa\_banach\_001, agent\_lane\_06}
\date{2026-06-28}

\begin{document}
\maketitle

\section*{Status}

This is a \texttt{literature\_already\_answered} packet.  It records that one
question from the final open-problems section of Alspach--Tong,
\emph{Subspaces of \(L_p\), \(p>2\), Determined by Partitions and Weights},
arXiv:math/0210228, was answered by the later Alspach--Tong note
\emph{Subspaces of \(L_p\), \(p>2\), with unconditional bases have equivalent
partition and weight norms}, arXiv:math/0310343.

\section*{Source Question}

In the ``Remarks and Open Problems'' section of arXiv:math/0210228,
Alspach and Tong ask whether \(L_p\) is isomorphic to a space whose norm is
given by partitions and weights.  This question occurs after their discussion
of the Bourgain--Rosenthal--Schechtman complemented subspaces and before their
question about constants in a Rosenthal-type inequality.  The source paper is
throughout concerned with the range \(2<p<\infty\).

\section*{Supporting Answer}

The abstract of arXiv:math/0310343 states that every subspace of \(L_p\),
\(2<p<\infty\), with an unconditional basis has an equivalent norm determined
by partitions and weights, and then explicitly says that consequently \(L_p\)
has such a norm.

The body of the paper gives the precise theorem:
\[
\text{if \(X\subset L_p\), \(2<p<\infty\), has an unconditional basis, then
\(X\) has an equivalent partition--weight norm.}
\]
The immediately following corollary states:
\[
L_p[0,1],\quad 2<p<\infty,
\text{ has an equivalent norm given by partitions and weights.}
\]
Since \(L_p[0,1]\) has an unconditional Haar basis for \(1<p<\infty\), this
corollary gives an affirmative answer, up to isomorphism, to the source
question in the \(p>2\) setting of arXiv:math/0210228.

\section*{Scope}

This packet is intentionally narrow.  It does not claim answers to the other
questions in the same open-problems section: existence of an envelope-property
example not embedding into \(L_p\), necessary and sufficient conditions for
partition--weight spaces to embed into \(L_p\), the tensor-product question,
the transfinite \(X_p^\alpha\) versus \(R_p^\alpha\) question, or sharp
constants in the displayed Rosenthal-type inequality.

\section*{Search Evidence}

The lane-6 sparse queue selected arXiv:math/0210228.  Cheap run indexes were
searched for the arXiv id, title words, \texttt{partitions and weights},
\texttt{envelope property}, \texttt{X\_p\^{}alpha}, \texttt{R\_p\^{}alpha},
and Rosenthal keywords; no duplicate packet for this source was found.  A
local full-source search for exact phrases from the open-problems section
found arXiv:math/0310343.  Its abstract, theorem, and corollary directly
identify the affirmative answer to the \(L_p\) question.

\begin{thebibliography}{9}
\bibitem{AT02}
D.~E. Alspach and S.~Tong,
\emph{Subspaces of \(L_p\), \(p>2\), Determined by Partitions and Weights},
arXiv:math/0210228.

\bibitem{AT03}
D.~E. Alspach and S.~Tong,
\emph{Subspaces of \(L_p\), \(p>2\), with unconditional bases have equivalent
partition and weight norms},
arXiv:math/0310343.
\end{thebibliography}

\end{document}
